\section{Contribution}

\begin{frame}{The Problem in the Running Example}
\begin{itemize}
    \item The manual BankSearch hierarchy is trusted, but its nodes are not intensional descriptions.
    \item The topic-model \ac{fca} hierarchy is explainable, but it may disagree with the manual hierarchy.
    \item A search system needs both: expert-compatible structure and document-based explanations.
\end{itemize}

\vspace{0.5em}
\begin{alertblock}{Core question}
Which BankSearch clusters are supported by both expert hierarchy constraints and document-derived topic structure?
\end{alertblock}
\end{frame}

\begin{frame}{Why Disagreement Is Expected}
\begin{columns}[T,onlytextwidth]
\begin{column}{0.48\textwidth}
\textbf{Manual hierarchy}
\begin{itemize}
    \item Encodes human category decisions.
    \item Uses topic labels and domain structure.
    \item Gives clean branches such as Sport $\rightarrow$ Soccer.
\end{itemize}
\end{column}
\begin{column}{0.48\textwidth}
\textbf{Topic-model FCA}
\begin{itemize}
    \item Encodes statistical regularities in text.
    \item Uses shared topic attributes.
    \item May mix categories when vocabulary overlaps.
\end{itemize}
\end{column}
\end{columns}

\vspace{0.6em}
\begin{block}{Consequence}
Exact agreement is a meaningful ideal, but a practical comparison must also detect partial semantic overlap.
\end{block}
\end{frame}






% \begin{frame}{State of the Art Around the Example}
% \begin{itemize}
%     \item Hierarchical text classification predicts labels inside a given taxonomy.
%     \item \ac{fca}-based hierarchy generation derives concepts directly from formal contexts.
%     \item Constrained clustering injects prior knowledge through constraints.
%     \item \ac{mlb} constraints express hierarchical merge order in agglomerative clustering.
% \end{itemize}

% \vspace{0.5em}
% \begin{block}{Gap}
% We have data-driven \ac{fca} hierarchies and constrained clustering, but no systematic way to use \ac{mlb} expert hierarchy constraints inside \ac{fca} and compare them with topic-model concepts.
% \end{block}
% \end{frame}





\begin{frame}{Our Answer for the BankSearch Story}
\begin{enumerate}
    \item Encode the manual BankSearch hierarchy as \ac{mlb} constraints.
    \item Prove that these constraints induce a valid \ac{fca} closure operator.
    \item Build an \ac{mlb}-constrained lattice as the expert-compatible cluster space.
    \item Compare it with the topic-model \ac{fca} lattice using consistency and coherence.
\end{enumerate}
\end{frame}

\begin{frame}{Central Idea}
\begin{center}
\begin{tikzpicture}[
    node distance=1.15cm,
    box/.style={draw,rounded corners=2pt,align=center,minimum width=3.0cm,minimum height=1.0cm},
    arr/.style={-{Latex[length=2mm]},thick}
]
\node[box] (manual) {Manual BankSearch\\hierarchy};
\node[box,right=of manual] (mlb) {\ac{mlb} constraints\\$\phi_{\Sigma}$};
\node[box,below=of manual] (docs) {BankSearch\\documents};
\node[box,right=of docs] (tm) {\ac{lda} topics\\$\phi_{TM}$};
\node[box,right=1.4cm of mlb,yshift=-0.58cm] (compare) {Consistency\\and coherence};
\draw[arr] (manual) -- (mlb);
\draw[arr] (docs) -- (tm);
\draw[arr] (mlb) -- (compare);
\draw[arr] (tm) -- (compare);
\end{tikzpicture}
\end{center}

\begin{block}{Interpretation}
The manual hierarchy becomes a formal closure system. The document evidence becomes another closure system. Their agreement can then be measured directly.
\end{block}
\end{frame}

\begin{frame}{Why We Need Two Comparison Notions}
\begin{columns}[T,onlytextwidth]
\begin{column}{0.48\textwidth}
\textbf{Consistency}
\begin{itemize}
    \item Strict exact agreement.
    \item A cluster must be closed in both lattices.
    \item Useful as a gold-standard diagnostic.
\end{itemize}
\end{column}
\begin{column}{0.48\textwidth}
\textbf{Coherence}
\begin{itemize}
    \item Relaxed semantic agreement.
    \item Clusters may overlap without being identical.
    \item Useful for heterogeneous knowledge sources.
\end{itemize}
\end{column}
\end{columns}

\vspace{0.5em}
\begin{block}{Role in the story}
If consistency is empty, that is not failure. It tells us the manual and data-driven hierarchies differ, and coherence shows where they still align.
\end{block}
\end{frame}
